% --- Archon physics formalization source begin ---
% archon:physics
% archon:covers PhyXMiniProblems/problem_phyx_mini_0163.lean
% archon:source-report reports/phyx_mini/problem_phyx_mini_0163.source.json

\chapter{Physics problem phyx\_mini\_0163}
\label{ch:PhyXMiniProblems_problem_phyx_mini_0163}

\paragraph{Problem source.}
\#\# Physical scenario

In figure, an isotropic point source of light \$S\$ is positioned at distance \$d\$ from a viewing screen \$A\$ and the light intensity \$I\_P\$ at point \$P\$ (level with \$S\$) is measured. Then a plane mirror \$M\$ is placed behind \$S\$ at distance \$d\$.

\#\# Question

By how much is \$I\_P\$ multiplied by the presence of the mirror?

\#\# Answer choices

- A: A: 0.95

- B: B: 1.0

- C: C: 1.05

- D: D: 1.11

\#\# Figure caption (auxiliary; use the image as primary evidence)

The image appears to be a simplified diagram or illustration involving mirrors and two points or objects labeled with letters. 

1. **Objects and Labels**:
   - There is a rectangular shape on the left labeled "M" which likely represents a mirror. The gradient shading suggests reflection.
   - On the right, there is another rectangular shape labeled "A", possibly an absorber or another type of surface.
   - A dot in the center is labeled "S".
   - Another dot on the right, near "A", is labeled "P".

2. **Spacing and Measurement**:
   - The diagram includes arrows pointing between "S" and the edges of "M" and "A", with the label "d" indicating distance. There are two of these distances marked with arrows: from "S" to the left (towards "M"), and from "S" to the right (towards "A").
  
3. **Relationships**:
   - The position of "S" is symmetrically between "M" and "A", suggesting a central or equidistant relation.
   - The element labeled "P" is placed between "S" and "A".

This diagram might illustrate a physical, optical, or conceptual setup involving reflections or projections, possibly related to physics or engineering.

\#\# Dataset metadata

Geometrical Optics; Physical Model Grounding Reasoning, Multi-Formula Reasoning

\paragraph{Recorded answer/context.}
D: D: 1.11

\paragraph{Figure/image path.}
/root/proposal\_for\_physic/science-mango/phyx\_mini\_run/phyx\_data/test\_image/163.png

\paragraph{Formalization target.}
create a compiling Lean file with sorry bodies at `PhyXMiniProblems/problem\_phyx\_mini\_0163.lean`. The Lean declarations must preserve the physical quantities, dimensions or dimensional roles, figure labels, governing-law hypotheses, and final relation expressed by this problem.
Use Mathlib/Physlib names found through LeanExplore where available. If a physics API is missing, introduce faithful local abstractions rather than scalar placeholder aliases.

\begin{theorem}[Physics formalization target]
\label{thm:physics:phyx_mini_0163:target}
\lean{PhyXMiniProblems.ProblemPhyXMini0163.problem_phyx_mini_0163}
\uses{def:physics:phyx-mini-0163:phyxminiproblems-problemphyxmini0163-irradianceinwattspersquaremetre, def:physics:phyx-mini-0163:phyxminiproblems-problemphyxmini0163-mirrorintensitysetup, def:physics:phyx-mini-0163:phyxminiproblems-problemphyxmini0163-matchestextualmirrorscreendiagram, def:physics:phyx-mini-0163:phyxminiproblems-problemphyxmini0163-hasphysicalmirrorintensityparameters, def:physics:phyx-mini-0163:phyxminiproblems-problemphyxmini0163-satisfiesidealplanemirrorgeometry, def:physics:phyx-mini-0163:phyxminiproblems-problemphyxmini0163-satisfiesidealpointsourcemirrorlaws, def:physics:phyx-mini-0163:phyxminiproblems-problemphyxmini0163-isuniquematchingdisplayedmultiplier}
The assigned autoformalize agent should translate this physics problem into Lean declarations in the covered file, with theorem and lemma proof bodies written as `by sorry`.
\end{theorem}
\begin{proof}
This is an autoformalization task, not a proof task. Produce statements that can later be proved by the physics prover without weakening the physical model.
\end{proof}

% --- Archon declaration topology begin ---
\section{Declaration topology}
\begin{definition}[Optical Plane]
  \label{def:physics:phyx-mini-0163:phyxminiproblems-problemphyxmini0163-opticalplane}
  \lean{PhyXMiniProblems.ProblemPhyXMini0163.OpticalPlane}
  The Euclidean cross-section containing the mirror, source, and screen
\end{definition}
\begin{proof}
  This is the declared model definition of Optical Plane.
\end{proof}
\begin{definition}[Length Quantity]
  \label{def:physics:phyx-mini-0163:phyxminiproblems-problemphyxmini0163-lengthquantity}
  \lean{PhyXMiniProblems.ProblemPhyXMini0163.LengthQuantity}
  A nonnegative physical length, independent of the chosen unit readout
\end{definition}
\begin{proof}
  This is the declared model definition of Length Quantity.
\end{proof}
\begin{definition}[Radiant Power Quantity]
  \label{def:physics:phyx-mini-0163:phyxminiproblems-problemphyxmini0163-radiantpowerquantity}
  \lean{PhyXMiniProblems.ProblemPhyXMini0163.RadiantPowerQuantity}
  A nonnegative radiant power, with physical dimension mass * length² / time³
\end{definition}
\begin{proof}
  This is the declared model definition of Radiant Power Quantity.
\end{proof}
\begin{definition}[Irradiance Quantity]
  \label{def:physics:phyx-mini-0163:phyxminiproblems-problemphyxmini0163-irradiancequantity}
  \lean{PhyXMiniProblems.ProblemPhyXMini0163.IrradianceQuantity}
  A nonnegative optical irradiance (power per area)
\end{definition}
\begin{proof}
  This is the declared model definition of Irradiance Quantity.
\end{proof}
\begin{definition}[length In Metres]
  \label{def:physics:phyx-mini-0163:phyxminiproblems-problemphyxmini0163-lengthinmetres}
  \lean{PhyXMiniProblems.ProblemPhyXMini0163.lengthInMetres}
  \uses{def:physics:phyx-mini-0163:phyxminiproblems-problemphyxmini0163-lengthquantity}
  Read a physical length as a scalar number of SI metres
\end{definition}
\begin{proof}
  This is the declared model definition of length In Metres.
\end{proof}
\begin{definition}[radiant Power In Watts]
  \label{def:physics:phyx-mini-0163:phyxminiproblems-problemphyxmini0163-radiantpowerinwatts}
  \lean{PhyXMiniProblems.ProblemPhyXMini0163.radiantPowerInWatts}
  \uses{def:physics:phyx-mini-0163:phyxminiproblems-problemphyxmini0163-radiantpowerquantity}
  Read a radiant power as a scalar number of SI watts
\end{definition}
\begin{proof}
  This is the declared model definition of radiant Power In Watts.
\end{proof}
\begin{definition}[irradiance In Watts Per Square Metre]
  \label{def:physics:phyx-mini-0163:phyxminiproblems-problemphyxmini0163-irradianceinwattspersquaremetre}
  \lean{PhyXMiniProblems.ProblemPhyXMini0163.irradianceInWattsPerSquareMetre}
  \uses{def:physics:phyx-mini-0163:phyxminiproblems-problemphyxmini0163-irradiancequantity}
  Read an irradiance as a scalar number of SI watts per square metre
\end{definition}
\begin{proof}
  This is the declared model definition of irradiance In Watts Per Square Metre.
\end{proof}
\begin{definition}[Isotropic Point Light Source]
  \label{def:physics:phyx-mini-0163:phyxminiproblems-problemphyxmini0163-isotropicpointlightsource}
  \lean{PhyXMiniProblems.ProblemPhyXMini0163.IsotropicPointLightSource}
  \uses{def:physics:phyx-mini-0163:phyxminiproblems-problemphyxmini0163-opticalplane, def:physics:phyx-mini-0163:phyxminiproblems-problemphyxmini0163-radiantpowerquantity}
  The isotropic point source S, including its position and radiant power
\end{definition}
\begin{proof}
  This is the declared model definition of Isotropic Point Light Source.
\end{proof}
\begin{definition}[Plane Mirror]
  \label{def:physics:phyx-mini-0163:phyxminiproblems-problemphyxmini0163-planemirror}
  \lean{PhyXMiniProblems.ProblemPhyXMini0163.PlaneMirror}
  \uses{def:physics:phyx-mini-0163:phyxminiproblems-problemphyxmini0163-opticalplane}
  An idealized plane mirror in the two-dimensional cross-section. Its affine line is the complete carrier used by the method of images. The reference point is the point on M level with S in the textual diagram
\end{definition}
\begin{proof}
  This is the declared model definition of Plane Mirror.
\end{proof}
\begin{definition}[virtual Image]
  \label{def:physics:phyx-mini-0163:phyxminiproblems-problemphyxmini0163-planemirror-virtualimage}
  \lean{PhyXMiniProblems.ProblemPhyXMini0163.PlaneMirror.virtualImage}
  \uses{def:physics:phyx-mini-0163:phyxminiproblems-problemphyxmini0163-opticalplane, def:physics:phyx-mini-0163:phyxminiproblems-problemphyxmini0163-planemirror}
  The virtual image of a point under reflection in the mirror carrier
\end{definition}
\begin{proof}
  This is the declared model definition of virtual Image.
\end{proof}
\begin{definition}[Viewing Screen]
  \label{def:physics:phyx-mini-0163:phyxminiproblems-problemphyxmini0163-viewingscreen}
  \lean{PhyXMiniProblems.ProblemPhyXMini0163.ViewingScreen}
  \uses{def:physics:phyx-mini-0163:phyxminiproblems-problemphyxmini0163-opticalplane}
  The viewing screen A and its marked observation point P
\end{definition}
\begin{proof}
  This is the declared model definition of Viewing Screen.
\end{proof}
\begin{definition}[Mirror Intensity Setup]
  \label{def:physics:phyx-mini-0163:phyxminiproblems-problemphyxmini0163-mirrorintensitysetup}
  \lean{PhyXMiniProblems.ProblemPhyXMini0163.MirrorIntensitySetup}
  \uses{def:physics:phyx-mini-0163:phyxminiproblems-problemphyxmini0163-opticalplane, def:physics:phyx-mini-0163:phyxminiproblems-problemphyxmini0163-lengthquantity, def:physics:phyx-mini-0163:phyxminiproblems-problemphyxmini0163-irradiancequantity, def:physics:phyx-mini-0163:phyxminiproblems-problemphyxmini0163-isotropicpointlightsource, def:physics:phyx-mini-0163:phyxminiproblems-problemphyxmini0163-planemirror, def:physics:phyx-mini-0163:phyxminiproblems-problemphyxmini0163-viewingscreen}
  All named quantities in the mirror/source/screen experiment. directAtP is the measured baseline irradiance before the mirror is installed; reflectedAtP is the mirror contribution; and totalWithMirrorAtP is the irradiance after installation. No field contains their requested ratio
\end{definition}
\begin{proof}
  This is the declared model definition of Mirror Intensity Setup.
\end{proof}
\begin{definition}[Matches Textual Mirror Screen Diagram]
  \label{def:physics:phyx-mini-0163:phyxminiproblems-problemphyxmini0163-matchestextualmirrorscreendiagram}
  \lean{PhyXMiniProblems.ProblemPhyXMini0163.MatchesTextualMirrorScreenDiagram}
  \uses{def:physics:phyx-mini-0163:phyxminiproblems-problemphyxmini0163-lengthinmetres, def:physics:phyx-mini-0163:phyxminiproblems-problemphyxmini0163-mirrorintensitysetup}
  The two distance labels and the horizontal alignment stated by the textual diagram. The unit optical-axis vector points from the mirror through S toward P; hence M--S and S--P each have length d
\end{definition}
\begin{proof}
  This is the declared model definition of Matches Textual Mirror Screen Diagram.
\end{proof}
\begin{definition}[Has Physical Mirror Intensity Parameters]
  \label{def:physics:phyx-mini-0163:phyxminiproblems-problemphyxmini0163-hasphysicalmirrorintensityparameters}
  \lean{PhyXMiniProblems.ProblemPhyXMini0163.HasPhysicalMirrorIntensityParameters}
  \uses{def:physics:phyx-mini-0163:phyxminiproblems-problemphyxmini0163-lengthinmetres, def:physics:phyx-mini-0163:phyxminiproblems-problemphyxmini0163-radiantpowerinwatts, def:physics:phyx-mini-0163:phyxminiproblems-problemphyxmini0163-irradianceinwattspersquaremetre, def:physics:phyx-mini-0163:phyxminiproblems-problemphyxmini0163-mirrorintensitysetup}
  Strictly positive distance, emitted power, and measured baseline irradiance. These exclude the degenerate cases in which an intensity multiplier is not defined
\end{definition}
\begin{proof}
  This is the declared model definition of Has Physical Mirror Intensity Parameters.
\end{proof}
\begin{definition}[Satisfies Ideal Plane Mirror Geometry]
  \label{def:physics:phyx-mini-0163:phyxminiproblems-problemphyxmini0163-satisfiesidealplanemirrorgeometry}
  \lean{PhyXMiniProblems.ProblemPhyXMini0163.SatisfiesIdealPlaneMirrorGeometry}
  \uses{def:physics:phyx-mini-0163:phyxminiproblems-problemphyxmini0163-planemirror-virtualimage, def:physics:phyx-mini-0163:phyxminiproblems-problemphyxmini0163-mirrorintensitysetup}
  The ideal plane-mirror image law along the normal optical axis. It is stated for every signed axial displacement, independently of the displayed value d and of any irradiance. Thus it specifies the mirror geometry without assuming the requested multiplier
\end{definition}
\begin{proof}
  This is the declared model definition of Satisfies Ideal Plane Mirror Geometry.
\end{proof}
\begin{definition}[Satisfies Inverse Square Irradiance]
  \label{def:physics:phyx-mini-0163:phyxminiproblems-problemphyxmini0163-satisfiesinversesquareirradiance}
  \lean{PhyXMiniProblems.ProblemPhyXMini0163.SatisfiesInverseSquareIrradiance}
  \uses{def:physics:phyx-mini-0163:phyxminiproblems-problemphyxmini0163-opticalplane, def:physics:phyx-mini-0163:phyxminiproblems-problemphyxmini0163-radiantpowerquantity, def:physics:phyx-mini-0163:phyxminiproblems-problemphyxmini0163-irradiancequantity, def:physics:phyx-mini-0163:phyxminiproblems-problemphyxmini0163-radiantpowerinwatts, def:physics:phyx-mini-0163:phyxminiproblems-problemphyxmini0163-irradianceinwattspersquaremetre}
  The inverse-square relation for one isotropic point emitter. powerFraction is dimensionless; it is 1 for the direct source and the mirror reflectance for the virtual source. This governing law contains no source/mirror distance ratio and no answer-choice value
\end{definition}
\begin{proof}
  This is the declared model definition of Satisfies Inverse Square Irradiance.
\end{proof}
\begin{definition}[Satisfies Ideal Point Source Mirror Laws]
  \label{def:physics:phyx-mini-0163:phyxminiproblems-problemphyxmini0163-satisfiesidealpointsourcemirrorlaws}
  \lean{PhyXMiniProblems.ProblemPhyXMini0163.SatisfiesIdealPointSourceMirrorLaws}
  \uses{def:physics:phyx-mini-0163:phyxminiproblems-problemphyxmini0163-irradianceinwattspersquaremetre, def:physics:phyx-mini-0163:phyxminiproblems-problemphyxmini0163-planemirror-virtualimage, def:physics:phyx-mini-0163:phyxminiproblems-problemphyxmini0163-mirrorintensitysetup, def:physics:phyx-mini-0163:phyxminiproblems-problemphyxmini0163-satisfiesinversesquareirradiance}
  The physical laws used in the computation: the installed mirror is perfectly reflecting; the direct and virtual-image contributions obey the inverse-square law; and geometrical-optics irradiances add incoherently at P
\end{definition}
\begin{proof}
  This is the declared model definition of Satisfies Ideal Point Source Mirror Laws.
\end{proof}
\begin{lemma}[virtual Image to P distance eq three d]
  \label{lem:physics:phyx-mini-0163:phyxminiproblems-problemphyxmini0163-virtualimage-to-p-distance-eq-three-d}
  \lean{PhyXMiniProblems.ProblemPhyXMini0163.virtualImage_to_P_distance_eq_three_d}
  \uses{def:physics:phyx-mini-0163:phyxminiproblems-problemphyxmini0163-lengthinmetres, def:physics:phyx-mini-0163:phyxminiproblems-problemphyxmini0163-planemirror-virtualimage, def:physics:phyx-mini-0163:phyxminiproblems-problemphyxmini0163-mirrorintensitysetup, def:physics:phyx-mini-0163:phyxminiproblems-problemphyxmini0163-matchestextualmirrorscreendiagram, def:physics:phyx-mini-0163:phyxminiproblems-problemphyxmini0163-hasphysicalmirrorintensityparameters, def:physics:phyx-mini-0163:phyxminiproblems-problemphyxmini0163-satisfiesidealplanemirrorgeometry}
  The mirror image of S is three marked distances from P
\end{lemma}
\begin{proof}
  The conclusion follows from the governing relations and figure readouts named in the statement of virtual Image to P distance eq three d.
\end{proof}
\begin{lemma}[reflected irradiance eq one ninth direct]
  \label{lem:physics:phyx-mini-0163:phyxminiproblems-problemphyxmini0163-reflected-irradiance-eq-one-ninth-direct}
  \lean{PhyXMiniProblems.ProblemPhyXMini0163.reflected_irradiance_eq_one_ninth_direct}
  \uses{def:physics:phyx-mini-0163:phyxminiproblems-problemphyxmini0163-irradianceinwattspersquaremetre, def:physics:phyx-mini-0163:phyxminiproblems-problemphyxmini0163-mirrorintensitysetup, def:physics:phyx-mini-0163:phyxminiproblems-problemphyxmini0163-matchestextualmirrorscreendiagram, def:physics:phyx-mini-0163:phyxminiproblems-problemphyxmini0163-hasphysicalmirrorintensityparameters, def:physics:phyx-mini-0163:phyxminiproblems-problemphyxmini0163-satisfiesidealplanemirrorgeometry, def:physics:phyx-mini-0163:phyxminiproblems-problemphyxmini0163-satisfiesidealpointsourcemirrorlaws}
  By the inverse-square law, the reflected contribution is 1/9 of baseline
\end{lemma}
\begin{proof}
  The conclusion follows from the governing relations and figure readouts named in the statement of reflected irradiance eq one ninth direct.
\end{proof}
\begin{definition}[Answer Choice]
  \label{def:physics:phyx-mini-0163:phyxminiproblems-problemphyxmini0163-answerchoice}
  \lean{PhyXMiniProblems.ProblemPhyXMini0163.AnswerChoice}
  Labels of the four displayed multiple-choice multipliers
\end{definition}
\begin{proof}
  This is the declared model definition of Answer Choice.
\end{proof}
\begin{definition}[displayed Multiplier]
  \label{def:physics:phyx-mini-0163:phyxminiproblems-problemphyxmini0163-displayedmultiplier}
  \lean{PhyXMiniProblems.ProblemPhyXMini0163.displayedMultiplier}
  \uses{def:physics:phyx-mini-0163:phyxminiproblems-problemphyxmini0163-answerchoice}
  The decimal multiplier printed beside each answer label
\end{definition}
\begin{proof}
  This is the declared model definition of displayed Multiplier.
\end{proof}
\begin{definition}[recorded Dataset Answer]
  \label{def:physics:phyx-mini-0163:phyxminiproblems-problemphyxmini0163-recordeddatasetanswer}
  \lean{PhyXMiniProblems.ProblemPhyXMini0163.recordedDatasetAnswer}
  \uses{def:physics:phyx-mini-0163:phyxminiproblems-problemphyxmini0163-answerchoice}
  The dataset's recorded label, retained as metadata rather than a premise
\end{definition}
\begin{proof}
  This is the declared model definition of recorded Dataset Answer.
\end{proof}
\begin{definition}[Matches Displayed Multiplier]
  \label{def:physics:phyx-mini-0163:phyxminiproblems-problemphyxmini0163-matchesdisplayedmultiplier}
  \lean{PhyXMiniProblems.ProblemPhyXMini0163.MatchesDisplayedMultiplier}
  \uses{def:physics:phyx-mini-0163:phyxminiproblems-problemphyxmini0163-irradianceinwattspersquaremetre, def:physics:phyx-mini-0163:phyxminiproblems-problemphyxmini0163-mirrorintensitysetup, def:physics:phyx-mini-0163:phyxminiproblems-problemphyxmini0163-answerchoice, def:physics:phyx-mini-0163:phyxminiproblems-problemphyxmini0163-displayedmultiplier}
  A choice matches the physical multiplier when its displayed two-decimal value is within half of 0.01 of the exact intensity ratio
\end{definition}
\begin{proof}
  This is the declared model definition of Matches Displayed Multiplier.
\end{proof}
\begin{definition}[Is Unique Matching Displayed Multiplier]
  \label{def:physics:phyx-mini-0163:phyxminiproblems-problemphyxmini0163-isuniquematchingdisplayedmultiplier}
  \lean{PhyXMiniProblems.ProblemPhyXMini0163.IsUniqueMatchingDisplayedMultiplier}
  \uses{def:physics:phyx-mini-0163:phyxminiproblems-problemphyxmini0163-mirrorintensitysetup, def:physics:phyx-mini-0163:phyxminiproblems-problemphyxmini0163-answerchoice, def:physics:phyx-mini-0163:phyxminiproblems-problemphyxmini0163-matchesdisplayedmultiplier}
  A displayed choice is the unique two-decimal match to the physical ratio
\end{definition}
\begin{proof}
  This is the declared model definition of Is Unique Matching Displayed Multiplier.
\end{proof}
% --- Archon declaration topology end ---
% --- Archon physics formalization source end ---
